\scanpage{218}
\mspage{167}

de $G(x)\setminus F(x)$, et à cause du caractère cocartésien de $Q$,

\[
(G(x)\mathrel{\unsure{-|\setminus}}F(x))\mathrel{\unsure{-|\longrightarrow}}(G'(x)\unsure{>|\setminus}F'(x))
\]

est bijectif, i.e. ``il revient au $=$'' de se donner $x\mathrel{\unsure{-|\longrightarrow}}G$ ne se factorisant pas par $F$, ou $\unsure{a|x}\longrightarrow G'$ ne se factorisant pas par $F'$.

Ainsi on a prouvé le

\uline{Théorème 5} : Soit $X$ dans Cat, et

\diagram{\begin{tikzcd}[ampersand replacement=\&]
Q\colon \& F\arrow[r,hook,no head,"i"]\arrow[d,"\unsure{p|r}"']\& G\arrow[d,"q"]\\
\& F'\arrow[r,"i'"]\& G'
\end{tikzcd}}

un carré cocartésien dans $X^\wedge$, avec $i$ mono. Alors $i'$ est mono, le carré \inserted{$Q$} est aussi cartésien, et le carré correspondant dans $\mathrm{Cat}_{/X}$ ou dans Cat

\diagram{\begin{tikzcd}[ampersand replacement=\&]
X/F\arrow[r,"X/i"]\arrow[d,"X/p"'] \& X/G\arrow[d,dash,"X/q"]\\
X/\unsure{p'|P'|F'}\arrow[r,"X/i'"] \& X/G'
\end{tikzcd}}
% archive sheet number: 218
